% Source: https://artofproblemsolving.com/community/c3h1368798_the_number_of_data_points__1__2
% Stable identity: aops-topic-1368798-q01
The median is the middle of the ordered list: as many values lie before it as after it.

If $n$ is odd, say $n=2k+1$, the value in position $k+1$ has $k$ values before it and $k$ after it, so it is the median, and $k+1=\dfrac{n+1}{2}$.

If $n$ is even, say $n=2k$, no single position is in the middle: positions $k$ and $k+1$ each have $k-1$ values on their outer side. The median is defined as the mean of these two middle values. The formula $\dfrac{n+1}{2}=k+\dfrac{1}{2}$ is then not a whole number; it says the median sits halfway between positions $k=\dfrac {n}{2}$ and $k+1=\dfrac {n}{2}+1$.

(a) Sorted: $1,3,4,7,9,12,15$. Here $n=7$ and $\dfrac{n+1}{2}=4$, so the median is the $4$th value, $\boxed{7}$.

(b) Sorted: $1,3,4,7,9,10,12,15$. Here $n=8$ and $\dfrac{n+1}{2}=4.5$, so the median is halfway between the $4$th and $5$th values: $\dfrac{7+9}{2}=\boxed{8}$.
